% from main.tex
%%%%%%%%%%  scan idx 138  |  folio 125  %%%%%%%%%%
un
triangle distingué de \(\Dpx(\mathcal C)\). On peut dresser un tableau analogue à (2.6), où les cases indiquent des amplitudes parfaites et chaque ligne correspond à une implication dont le but est souligné.
\[
\begin{array}{ccc}
E&F&G\\[0.3em]
{}[a+1,b+1]&{}[a,b]&\underline{[a,b]}\\
{}[a,b]&\underline{[a,b]}&{}[a,b]\\
\underline{[a,b]}&{}[a,b]&{}[a-1,b-1].
\end{array}
\]
De (4.10) on déduit aussitôt le

\textbf{Corollaire 4.12.} \emph{Supposons que \(S\) soit un site ponctuel. Pour que \(F\in\Ob\Dpx(\mathcal C)\) soit d'amplitude parfaite (rel.\ à \(\mathcal C_0\)) contenue dans \([a,b]\), il faut et il suffit qu'il existe un quasi-isomorphisme \(F'\to F\), avec \(F'\) strictement parfait et tel que \(F'^i=0\) pour \(i\notin[a,b]\). Les foncteurs canoniques}
\[
  \Kpx^b(\mathcal C_0)\longrightarrow\Parf(\mathcal C,\mathcal C_0)
\]
\emph{et}
\[
  \Tr(\Kpx^b(\mathcal C_0))\longrightarrow\Tr(\Parf(\mathcal C,\mathcal C_0))
\]
\emph{sont des équivalences de catégories}.

\textbf{Remarque 4.12.1.} Quand \(S\) est un site ponctuel, les objets de \(\mathcal C_0\) sont projectifs. Lorsque \(\mathcal C_0\) contient tous les objets projectifs de \(\mathcal C\), on dit « amplitude projective » au lieu de « amplitude \(\mathcal C_0\)-parfaite ».



% Additional notation for Expose I continuation
\providecommand{\Cpx}{\mathrm C}
\providecommand{\Kpx}{\mathrm K}
\providecommand{\Dpx}{\mathrm D}
\providecommand{\Ho}{\mathrm H}
\providecommand{\RHom}{\mathrm R\!\operatorname{Hom}}
\renewcommand{\Ob}{\operatorname{ob}}
\providecommand{\qis}{\xrightarrow{\sim}}
\providecommand{\RHom}{\mathrm R\!\Hom}
\providecommand{\Ext}{\operatorname{Ext}}
\providecommand{\Hom}{\operatorname{Hom}}
\providecommand{\Tr}{\operatorname{Tr}}

\providecommand{\Add}{\operatorname{Add}}
\providecommand{\Cart}{\operatorname{Cart}}
\providecommand{\perfamp}{\operatorname{perf.amp}}
\providecommand{\perfdim}{\operatorname{perf.dim}}
\providecommand{\toramp}{\operatorname{tor.amp}}
\providecommand{\torf}{\operatorname{torf}}
\providecommand{\injf}{\operatorname{injf}}
\providecommand{\rg}{\operatorname{rg}}
\providecommand{\is}{\mathrm{is}}
\providecommand{\Mod}{\operatorname{Mod}}


